% =========================================================================
% Preface
% =========================================================================
\chapter*{Preface}
\addcontentsline{toc}{chapter}{Preface}
\markboth{Preface}{}

Poraqu\^e is a research code for learning maps between the three-dimensional
scalar fields of density-functional theory. It has one organising idea: for a
given material, the local external potential $\vext(\rr)$, the valence charge
density $\rho(\rr)$ and the kinetic energy density $\tau(\rr)$ all live on
\emph{one} grid, in \emph{one} file format, and are therefore directly
comparable, composable, and usable as aligned inputs and targets for a neural
operator.

Two maps are learned:
\[
  \vext \;\longmapsto\; \rho
  \qquad\text{and}\qquad
  \rho \;\longmapsto\; \tau .
\]
They are not two unrelated regressions that happen to share a file format. They
are the two legs of one variational statement --- the first is the
Hohenberg--Kohn map, guaranteed to exist by theorem; the second is the kinetic
energy density functional, the missing ingredient of orbital-free DFT. Composed,
they constitute a complete orbital-free calculation. That relationship is the
thread running through this document.

\section*{What this guide covers}

\begin{description}
  \item[Chapters 1--2] The physics. Kohn--Sham DFT, the Hohenberg--Kohn
    theorems, and the orbital-free formulation, with attention to which
    quantities are exact, which are approximations, and which are the open
    problem.
  \item[Chapter 3] The mathematics of Fourier neural operators: why an operator
    between function spaces is the right object, how the spectral convolution
    is constructed, and why it is the natural architecture for a periodic
    crystal on a plane-wave grid.
  \item[Chapters 4--5] Implementation. The shared-grid data model, the
    code-agnostic ingestion layer, and a detailed account of the interface to
    VASP --- including the exact reconstruction of the local pseudopotential.
  \item[Chapter 6] Measured results on a gold data set, with the analytic
    baselines a learned functional must beat.
  \item[Chapter 7] The physics-informed programme: how the governing equations
    re-enter as constraints, and which of them are already enforced by
    construction.
\end{description}

\section*{Conventions}

Two unit systems appear, deliberately.

\begin{itemize}
  \item \textbf{Å and eV} in everything that touches a file. Every plane-wave
    code writes geometry in Ångström and energies in electronvolt, so the field
    layer, the readers and the machine-learning code use those units
    throughout. A charge density is in $e/\text{\AA}^3$, a kinetic energy
    density in $\mathrm{eV}/\text{\AA}^3$, a potential in eV.
  \item \textbf{Hartree atomic units} in the analytical expressions of
    Chapters~1--2 and in the legacy solver, where $\hbar = m_e = e = 1$ removes
    constants that would otherwise clutter every equation.
\end{itemize}

Conversions are collected in \file{poraque/fields/constants.py}. The two
constants that recur are
\[
  \frac{e^2}{4\pi\varepsilon_0} = 14.399645\ \mathrm{eV}\cdot\text{\AA},
  \qquad
  \frac{\hbar^2}{2m_e} = 3.809982\ \mathrm{eV}\cdot\text{\AA}^2 .
\]

Sign convention for potentials: $\vext(\rr)$ is the potential \emph{energy of an
electron}, so it is negative near the ions. This matches VASP's \file{LOCPOT}
and \file{EXTCAR}.

\begin{pnote}
Passages marked \textbf{Implementation} connect an equation to the function
that evaluates it. Passages marked \textbf{Caveat} record an approximation, an
untested path, or a limit on what the numbers presented here support. They are
not disclaimers added for politeness: each one marks a place where a reader
could otherwise draw a conclusion the evidence does not carry.
\end{pnote}
